\chapter{Closures, Lambda Lifting, and Combinators}
\label{ch:closures-lifting-combinators}

A nested function that escapes its defining call cannot keep using a stack slot that has already disappeared.
Its free variables must travel with it or become explicit parameters.
Closure conversion, lambda lifting, and combinator translation solve this problem with different intermediate representations and costs.
They should be compared by allocation, access, calling behavior, and preserved evaluation semantics rather than by syntax alone.

\section{Compute free variables by binding identity}
\label{sec:ch91-free}
The free variables of a lambda body exclude its bound parameter and include references resolved outside the lambda.
For lambda y.add(x,y), x is free and y is bound.
For nested lets, the value expression of a nonrecursive let does not see its new binder, while the body does.
A recursive let includes the binder in the scope of its right side as well.

Algorithm~\ref{alg:ch91-free} collects free declaration identities.
With explicit set unions, cost depends on set sizes and repeated copying; a single traversal with scope tracking can avoid naive quadratic accumulation in long nested syntax.
Names must be resolved first or alpha-renaming must ensure unambiguous binders.
Using spelling alone after shadowing can capture the wrong x even if the resulting closure layout looks plausible.

\begin{algorithm}[htbp]
\caption{Collect the free declaration identities needed by a closure.}
\label{alg:ch91-free}
\begin{pseudocode}
\Require resolved syntax and parameter/local binding identities
\Ensure exactly the references supplied by the enclosing environment
\State traverse the body, maintaining the set of locally bound declarations
\ForAll{variable references}
\If{reference is not locally bound} \State add its declaration identity to captures \EndIf
\EndFor
\State include bindings needed by nested closures that are constructed here
\State return captures in a deterministic layout order
\end{pseudocode}
\end{algorithm}

The last step matters when an outer function does not itself read x but constructs an inner closure that does.
The outer code must retain enough information to populate the inner environment.
Deterministic field ordering simplifies generated-code tests and debug metadata.
It also creates an ABI-like contract between closure construction and the compiled closure body.

\section{Convert closures into code and environment}
\label{sec:ch91-convert}
Closure conversion turns the function into a code label plus an environment record.
The converted code takes an explicit environment argument and accesses captured fields through known offsets.
Calling a closure loads its code and environment and supplies the ordinary argument under the chosen calling convention.
Figure~\ref{fig:ch91-layout} shows a closure capturing x=11.

\begin{figure}[htbp]
\centering
\begin{tikzpicture}[node distance=15mm and 18mm]
\node[basic block] (closure) {closure: code, environment};
\node[basic block,below left=of closure] (code) {code(env,y): env.x + y};
\node[basic block,below right=of closure] (env) {field x = 11};
\draw[cfg edge] (closure)--(code);\draw[cfg edge] (closure)--(env);
\end{tikzpicture}
\caption{Converted code and environment must agree on capture layout and calling convention.}
\label{fig:ch91-layout}
\end{figure}

A flat closure copies needed values into one record; a linked environment retains parent frames and follows access paths.
Flat access is simple but can duplicate captures across closures.
Linked access can share frames but retain unrelated values and add indirections.
Capturing mutable variables requires retaining the shared cell, not copying its current value, if later mutations must be visible to every closure.
The pure lab uses cells primarily for recursive bindings and lexical environments.

Table~\ref{tab:ch91-layouts} lists these tradeoffs without assigning universal performance winners.
Escape analysis can keep some environments on the stack or eliminate them after inlining, but every escaping reference needs a valid lifetime.
The collector must trace captured references according to the environment layout.
A code pointer alone is not a complete function value when free variables remain.

\begin{table}[htbp]
\caption{Environment representation changes access cost, sharing, and retained memory.}
\label{tab:ch91-layouts}
\centering
\begin{tabular}{lll}
\toprule
Representation & Benefit & Cost or obligation \\
\midrule
flat environment & direct capture offsets & copied fields; layout metadata \\
linked frames & shared enclosing storage & access paths; accidental retention \\
captured cell & shared update identity & extra indirection and traced lifetime \\
eliminated closure & no indirect closure call & requires escape/inlining proof \\
\bottomrule
\end{tabular}
\end{table}

\section{Lift free variables into explicit parameters}
\label{sec:ch91-lifting}
\term{Lambda lifting} moves suitable local functions to a wider scope and adds parameters for their free variables.
The local lambda y.add(x,y) can become a top-level function addCaptured(x,y).
Every transformed call supplies the correct x from its lexical environment.
If the function itself escapes as a value, the remaining partial application still needs a representation for the supplied x; lifting does not universally eliminate closure allocation.

Mutually recursive local functions require computing capture requirements across their call graph.
A function may need an extra parameter solely to pass it to another lifted function.
Fixed-point propagation over call dependencies can determine those requirements, but parameter explosion and repeated argument passing can make unrestricted lifting unattractive.
Selective lifting measures the tradeoff rather than assuming fewer nested lambdas means faster code.

Algorithm~\ref{alg:ch91-lift} describes a restricted direct-call lifting transformation.
It assumes resolved direct callees, no escaping function value for the lifted group, and a fixed evaluation-order contract for capture expressions.
With $F$ functions and $C$ call edges, repeated set propagation costs the number of iterations times edge/set work; strongly connected components can organize recursive groups efficiently.
The lab explains this transformation but does not implement a production lifting pass.

\begin{algorithm}[htbp]
\caption{Lift a nonescaping direct-call group with propagated capture requirements.}
\label{alg:ch91-lift}
\begin{pseudocode}
\Require resolved local function group, known direct calls, no escaping function values
\Ensure lifted definitions and calls preserving lexical bindings
\State collect each body's immediate free-variable requirements
\State propagate callee capture requirements to callers until stable
\State choose deterministic added-parameter order for each lifted function
\State move definitions and rewrite captured references to added parameters
\State rewrite every call to supply its callee's required lexical values
\State verify no former local binding remains free
\end{pseudocode}
\end{algorithm}

\section{Eliminate variables through combinators}
\label{sec:ch91-combinators}
Combinator translation replaces abstractions with closed operators such as I, K, and S.
Their characteristic reductions are $I\,x\to x$, $K\,x\,y\to x$, and $S\,f\,g\,x\to f\,x\,(g\,x)$.
For a bracket abstraction over x, x maps to I, a term not containing x can map to K applied to that term, and an application can be decomposed through S.
These rules define a translation idea, not a guarantee of compact code.

For lambda x.x the result is I.
For lambda x.c where c is independent of x, the result is K c.
For lambda x.f x, a direct S/K/I translation produces S(K f)I; an eta-like contraction to f needs its own semantic conditions under the chosen evaluation strategy.
Table~\ref{tab:ch91-combinators} records the simple cases.

\begin{table}[htbp]
\caption{Bracket abstraction removes a binder by making application structure explicit.}
\label{tab:ch91-combinators}
\centering
\begin{tabular}{ll}
\toprule
Abstraction & Elementary translation \\
\midrule
lambda x.x & I \\
lambda x.c, with x absent & K c \\
lambda x.f x, with x absent from f & S(K f)I \\
\bottomrule
\end{tabular}
\end{table}

Naive translation can expand code substantially and create many runtime application nodes.
Supercombinators and richer combinator sets improve granularity by compiling larger closed functions and handling their arguments together.
Graph reduction can share arguments rather than duplicate their full subgraphs.
But eliminating source variables does not eliminate the need to specify evaluation strategy, heap updates, or garbage collection.

Figure~\ref{fig:ch91-choices} shows three alternative paths from nested syntax.
They are not three obligatory sequential stages of the case-study compiler.
The implemented compiler uses ordinary closures and structured stack code; the other paths explain design choices in functional-language implementation literature.

\begin{figure}[htbp]
\centering
\begin{tikzpicture}[node distance=13mm and 12mm]
\node[basic block] (source) {nested functional syntax};
\node[basic block,below left=of source] (closure) {closures};
\node[basic block,below=of source] (lift) {lifted functions};
\node[basic block,below right=of source] (comb) {combinators};
\draw[cfg edge] (source)--(closure);\draw[cfg edge] (source)--(lift);\draw[cfg edge] (source)--(comb);
\end{tikzpicture}
\caption{Closure conversion, lifting, and combinator translation offer different representation tradeoffs.}
\label{fig:ch91-choices}
\end{figure}

\section{Optimize higher-order calls with justified target information}
\label{sec:ch91-analysis}
A known closure target permits direct-call specialization and sometimes environment elimination.
Higher-order control-flow analysis approximates which closures can reach each call site.
A singleton target set can justify devirtualization only when the approximation is sound and calling/environment conventions remain compatible.
An empty set caused by an unfinished analysis is not a proof that the call is unreachable.

Shivers develops higher-order control-flow recovery and its optimization uses~\cite{shivers-1991-control}.
The cost and precision depend on context abstraction and representation, not merely on using a graph worklist.
The lab does not implement that analysis; its lexical-scope and reentrant-closure examples instead test the baseline closure semantics against the stack compiler.
These are the semantics an advanced analysis must preserve.

\section{Propagate captures through a recursive call group}
\label{sec:ch91-transitive}
Suppose local function f directly calls local function g, and only g reads an enclosing variable x.
If g is lifted to require x explicitly, f must be able to supply x at its call even though f's own arithmetic never inspected it.
Thus f's lifting requirement contains more than its immediate free-variable reads.
If f and g call one another, propagating these requirements must reach a fixed point over the recursive group.

With a finite set of enclosing declarations, requirements grow monotonically by set inclusion and cannot grow forever.
Each newly discovered required declaration must eventually propagate to every relevant predecessor call edge.
A worklist indexed by callers avoids repeatedly scanning unrelated functions, while a strongly connected component groups mutually recursive propagation.
The actual cost includes set representation and the number of distinct function/declaration membership additions.
This finite dataflow argument is much clearer than describing lifting as merely adding whatever names occur in each printed body.

The optimization tradeoff remains after the fixed point is correct.
Passing x through a deep chain can add repeated call arguments; capturing one environment pointer may be cheaper on a particular target.
Inlining g into f can remove one call but duplicate code.
A selective lifting pass therefore separates the mandatory capture analysis from the profitability decision about which functions to lift.
The same analysis facts can be useful for closure trimming even when no function is moved to top level.

\section{Exercises}
\label{sec:ch91-exercises}
\begin{enumerate}
\item\label{ex:ch91-captures} List the captures of lambda y.add(x,y) and give a flat environment layout.
\item\label{ex:ch91-lift} Lift that lambda for known direct calls, naming the added parameter and rewritten call arguments.
\item\label{ex:ch91-comb} Reduce S(K f)I applied to a symbol a under the stated combinator rules.
\item\label{ex:ch91-escape} Explain why returning a lifted function partially applied to captured x can still require a closure-like value.
\end{enumerate}

\section{Further reading}
\label{sec:ch91-reading}
Appel treats closure conversion and environment representation~\cite{appel-1998-modern}.
Shivers explains higher-order flow information~\cite{shivers-1991-control}; the STG paper develops a contrasting non-strict closure and entry design~\cite{peytonjones-1992-stg}.
